\clearpage
\VerdierPhysicalPageAnchor{0241}%
\begingroup\small\setlength{\parskip}{0pt}\setlength{\emergencystretch}{1.5em}
\setlength{\abovedisplayskip}{4pt plus 2pt minus 2pt}%
\setlength{\belowdisplayskip}{4pt plus 2pt minus 2pt}%

\VerdierSemanticAnchor{ch4d-4-5-2}
\noindent\textbf{4.5.2.}\quad
Soient $F\colon D(\mathcal A)\longrightarrow\mathcal B$ un foncteur
cohomologique stationnaire à droite, $p$ un entier, $X$ un objet de
$\operatorname{comp}(\mathcal A)$ et :

\VerdierSemanticAnchor{ch4d-4-5-2-second-filtration}
\[
0=X(-\infty)\longrightarrow\cdots\longrightarrow X(q)
\longrightarrow X(q+1)\longrightarrow\cdots
\longrightarrow X(+\infty)=X
\]
la deuxième filtration canonique (4.3.1) de $X$. Le système inductif
$F^p(X(q))$ est stationnaire, i.e.\ il existe un entier $q_0$ tel que pour
tout $q$, $q_0\leq q\leq+\infty$, le morphisme :
\[
F^p(X(q))\longrightarrow F^p(X)
\]
soit un isomorphisme. En effet, avec les notations de (4.2.3), on a une
suite exacte :
\[
F^{p-1}\bigl(X(q,+\infty)\bigr)\longrightarrow F^p(X(q))
\longrightarrow F^p(X)\longrightarrow F^p\bigl(X(q,+\infty)\bigr),
\]
et le foncteur $F$ étant stationnaire à droite, il existe un entier $q_0$
tel que pour tout $q>q_0$ les objets $F^p(X(q,+\infty))$ et
$F^{p-1}(X(q,+\infty))$ soient nuls.

\VerdierSemanticAnchor{ch4d-4-5-2-first-filtration}
Soient de même $F\colon D(\mathcal A)\longrightarrow\mathcal B$ un
foncteur cohomologique stationnaire à droite, $p$ un entier, $Y$ un objet de
$\operatorname{comp}(\mathcal A)$ et :
\[
0=Y(-\infty)\longrightarrow\cdots\longrightarrow Y(q)
\longrightarrow Y(q+1)\longrightarrow\cdots
\longrightarrow Y(+\infty)=Y,
\]
la première filtration de $Y$ (4.1.2). Il existe un entier $q_0$ tel que
pour tout $q<q_0$ l'objet $F^p(Y(q))$ soit nul. La vérification analogue à
la précédente est laissée au lecteur qui pourra de même, en passant aux
catégories opposées, étudier les propriétés analogues des foncteurs
stationnaires à gauche.

\VerdierSemanticAnchor{ch4d-4-5-3}
\medskip\noindent\textbf{Proposition 4.5.3.}\quad

\VerdierSemanticAnchor{ch4d-4-5-3-a}
\noindent\textit{a) Soit $Y$ un objet de
$\operatorname{comp}^{*}(\mathcal A)$ (1.1.1)
($*=+,-,b,{}$ ``vide''). Soit
$F\colon D^*(\mathcal A)\longrightarrow\mathcal B$ un foncteur
cohomologique. La première suite spectrale du foncteur $F$ relative au
complexe $Y$ (4.4.1) dont les termes initiaux sont (4.4.2) :}
\[
\textit{$\mathrm{I}_1^{p,q}=F^q(Y^p)$}
\]
\textit{est stationnaire (chap.~II, 4.4.2) dans les cas suivants :}

\begin{enumerate}
\item[1)] \textit{$Y$ est un objet de $\operatorname{comp}^{b}(\mathcal A)$.}
\item[2)] \textit{$Y$ est un objet de $\operatorname{comp}^{+}(\mathcal A)$
et $F$ est stationnaire à droite.}
\item[3)] \textit{$Y$ est un objet de $\operatorname{comp}^{-}(\mathcal A)$
et $F$ est stationnaire à gauche.}
\item[4)] \textit{Le foncteur $F$ est stationnaire.}
\end{enumerate}

\VerdierSemanticAnchor{ch4d-4-5-3-b}
\noindent\textit{b) Soient $X$ un objet de $D^*(\mathcal A)$ (1.2.4) et
$F\colon D^*(\mathcal A)\longrightarrow\mathcal B$ un foncteur
cohomologique. La deuxième suite spectrale du foncteur $F$ relative au
complexe $X$ (4.4.5) est stationnaire dans les cas suivants :}
\endgroup
